\section{Method}

The residual method has four steps. First, choose a baseline population: prior
elections for the same unit, comparable reporting units in the current election,
neighboring units, demographic or registration covariates, or a combination of
those sources. Second, fit a baseline model that predicts turnout or candidate
share for each unit. Third, compute the residual
\[
  r_i = y_i - \widehat{y}_i,
\]
where $y_i$ is the observed turnout rate or vote share for unit $i$ and
$\widehat{y}_i$ is the model prediction. Fourth, standardize the residuals and
rank units by magnitude.

The statistic can be reported as a standardized residual, a studentized
residual, a robust score, or a rank within the baseline population. A simple
implementation might divide by an estimated residual standard deviation; a more
careful one can use robust regression, cross-validation, heavy-tailed errors,
or leave-one-out studentization so that one extreme unit does not define the
baseline that judges it. The report should identify the model as a residual
model and name the feature set rather than present the score as self-evident.

A unit can be a residual outlier for innocent reasons. New housing, a college
dormitory, a changed precinct boundary, a local uncontested race, language
access changes, weather, ballot style, or a mobilization campaign can all move
turnout or vote share away from what history predicted. The method therefore
ranks units for review; it does not explain the residual by itself.
